%==============================================================================
% Chapter 4 -- Results and discussion
%==============================================================================
\chapter{Results and discussion}
\label{ch:results}

We now show that Design-CP makes it possible to design whole symmetric
nanoparticles end to end, characterise how it scales with the number of GPUs, and
demonstrate that, because of the symmetry of the designed system, the approach is feasible even on
commodity hardware. We close with a discussion of its limitations.

\section{Design quality and symmetry}
\label{sec:design-quality}

Context parallelism is applied here as an \emph{inference-only} modification: it
changes how intermediate activations are partitioned and communicated, but does
not alter the learned parameters of RFD3. This makes it immediately applicable to
pretrained checkpoints, but it also means that when context parallelism is used to
exceed RFD3's native crop limits, the model is sampled outside the regime it was
trained on: $384$ tokens and $5{,}000$ atoms, according to the supplementary
material of \textcite{butcher_novo_2025}. In practice, we find that this
distribution shift can degrade sample quality when the target displays limited
symmetry.

Figure~\ref{fig:large-designs}b illustrates the effect qualitatively: when the
number of atoms and tokens modelled is much higher than RFD3 saw during training,
in this example a single monomer of $10{,}800$ amino acids sampled with 2D
sharding, we observe visibly unnatural designs. By contrast, when constraints are
introduced through symmetry while keeping the overall system size constant, the resulting
assemblies appear visibly more protein-like even while the system size is kept constant. 

We also assess that this apparent increase in sample quality is not an effect of chain length, 
by designing and visually checking a series of samples at exactly the icosahedral system size
($60$ chains of $210$ residues, $12{,}600$ residues in total) with \emph{no}
point-group constraint imposed at sampling time, keeping tokens, atoms,
architecture, weights and parallelisation strategy identical to the symmetric
runs. The resulting structures are visibly degenerate: large slabs of
$\beta$-strands packed in irregular orientations, no recognisable globular folds
and no consistent inter-chain organisation (Appendix~\ref{app:monstromers},
Figure~\ref{fig:monstromers}). We take this as direct qualitative evidence that
the dominant factor is the symmetry prior rather than the length of the modelled
chains.

We attribute this effect to the symmetry sampling mechanism described in
Section~\ref{sec:rfd3-prelim}: at each symmetrised denoising step the network
performs a full forward pass over all atoms of the assembly, but only the
asymmetric subunit (ASU) of its prediction is retained, and the remaining subunits
are overwritten by deterministic group-operation copies of that ASU. The sampler
therefore varies the coordinates of a single ASU rather than those of the full
assembly, which shrinks the effective degrees of freedom in the design space and
couples distant regions of the assembly by forcing every subunit to share the same
ASU backbone.

This motivates the use of context parallelism for the design of highly symmetric
protein assemblies such as octahedral and icosahedral capsids and cages. Context
parallelism allows the model to consider, and remain self-consistent with respect
to, the full set of residue and atom interactions in the assembly, while the
symmetry prior ensures that the number of free coordinates the sampler must
generate is that of a single ASU, allowing the use of pre-trained weights.
We turn to these targets next, as they are also our interest for nanoparticle vaccine design.

\section{Designing icosahedral nanoparticles}
\label{sec:icosahedral}

We generated $n{=}40$ icosahedral nanoparticles ($210$ residues per chain, $60$
chains, $12{,}600$ residues per assembly) with the 2D sharding scheme on a
$2\times2$ grid of NVIDIA GH200 GPUs ($95\,$GB each) at batch size one. As set out 
in Section~\ref{sec:evaluation}, we ask two questions:(i)~whether sampling at this
scale via Design-CP preserves the per-chain backbone quality of the original
RFD3 model, and (ii)~whether the symmetrised assembly realises the intended
interfaces, consistently with a functional natural baseline. As a paired control
for question~(i), we additionally sampled $n{=}40$ single-chain $210$-residue
monomers with the standard single-GPU RFD3. Symmetry-aware metrics for
question~(ii) are computed between the ASU and its eight nearest icosahedral
neighbours (Figure~\ref{fig:icosahedral}a), and the corresponding lumazine
synthase values are overlaid as a dotted reference line in every panel.


\begin{figure}[t]
  \centering
  \includegraphics[width=\textwidth]{Figure_2.png}
  \caption{\textbf{Designing icosahedral nanoparticles with Design-CP}
    ($n{=}40$; $60$ chains $\times\,210$ residues $=12{,}600$ residues).
    \textbf{(a)} A single asymmetric subunit (ASU, red) and its neighbouring
    subunits (teal) on the icosahedral shell.
    \textbf{(b--e)} Per-chain backbone-sanity metrics for the icosahedral designs
    (I60) versus $n{=}40$ single-GPU RFD3 monomers, with the lumazine-synthase
    reference (PDB~1NQX, dotted).
    \textbf{(f)} A selected design with a zoom on an inter-ASU interface annotated
    with C$\alpha$--C$\alpha$ distances.
    \textbf{(g--l)} Symmetry-aware interface metrics against the 1NQX baseline.
    Full distributions are reported in Appendix~\ref{app:extended}
    (Figures~\ref{fig:ico-perchain}--\ref{fig:ico-symm}).}
  \label{fig:icosahedral}
\end{figure}

\paragraph{Per-chain backbone quality broadly matches the original RFD3.}
On the backbone-sanity indicators (Figure~\ref{fig:icosahedral}b--e), the two
populations have broadly similar distributions: the average number of chain breaks
per chain and the average count of backbone heavy-atom clashes per chain both have
medians at zero. The ASUs extracted from the nanoparticles designed with Design-CP
show, however, more outliers than the designed monomers, especially for the number
of backbone clashes. The fraction of residues assigned to a helix or
$\beta$-strand element is closely matched between them and to the 1NQX reference
(medians of $\approx 0.60$ for the icosahedral ASUs against $\approx 0.65$ for
monomers, and $\approx 0.68$ for 1NQX). A more notable gap appears in the maximum
per-chain deviation of consecutive C$\alpha$--C$\alpha$ bond lengths from the
canonical \SI{3.8}{\angstrom} value: monomers have a distribution concentrated
around $\approx \SI{0.15}{\angstrom}$, whereas the icosahedral ASUs show a broader
bulk around $\approx \SI{0.30}{\angstrom}$. We partially attribute these
differences to the geometric strain of having to remain self-consistent with
$59$ symmetry mates and all of the inter-subunit interfaces they form, inside a
$12{,}600$-residue joint context, rather than to a degradation introduced by the
context-parallel sampler itself. The icosahedral design problem is strictly harder
per chain than designing monomers of the same length.

\paragraph{Symmetric interfaces track a functional natural assembly.}
Symmetry-aware metrics (Figure~\ref{fig:icosahedral}g--l) generally follow the
1NQX baseline, but not every sample produces a sterically clean assembly. The
number of inter-subunit clashes within the ASU's neighbourhood is non-zero for a
substantial fraction of designs, despite the mean being around zero.
Concordantly, the minimum C$\alpha$--C$\alpha$ distance between distinct chains is
centred near the 1NQX reference at $\approx \SI{4}{\angstrom}$, but the lower tail
of the distribution descends toward and below the \SI{3.5}{\angstrom}
steric-overlap threshold, indicating that a fraction of designs realise inter-chain
contacts inside the steric-overlap regime. The average number of
C$\alpha$--C$\alpha$ contacts per inter-subunit interface centres around
$\approx 125$, slightly above 1NQX rather than collapsing toward zero, and the
number of interfaces making contact resembles that of 1NQX. Combined with the
local visual evidence from the selected sample in
Figure~\ref{fig:icosahedral}f, where ASU--ASU contacts settle into well-formed
packing geometries with C$\alpha$--C$\alpha$ distances in the expected
\SIrange{4}{10}{\angstrom} band, this indicates that the symmetrised denoising
trajectory can recover physically reasonable inter-subunit interfaces rather than
independently designing non-interacting chains.

Together, these results show that Design-CP scales RFD3 sampling to large
icosahedral assemblies without retraining or fine-tuning, broadly preserving
per-chain backbone quality on par with vanilla RFD3 monomers and producing
symmetry-consistent interfaces comparable to a functional natural icosahedral
particle. Full metric distributions are reported in
Appendix~\ref{app:extended}.


\begin{table*}[!t]
  \centering
  \small
  \caption{\textbf{Scaling of Design-CP under icosahedral symmetry.} For each GPU count $P$,
   we report the maximum ASU length before out-of-memory and the wall-clock inference time 
   \emph{at that maximum ASU length}, for the 1D and 2D sharding schemes. 
   \emph{Capacity ratio} is defined as $L_{\max}^{\mathrm{2D}}/L_{\max}^{\mathrm{1D}}$ 
   and \emph{speedup} as $t_{\mathrm{1D}}/t_{\mathrm{2D}}$, so that values greater than 
   unity indicate that the 2D scheme outperforms the 1D scheme. All measurements use HG200 
   GPUs (95\,GB each); a dash denotes a configuration that cannot be measured.}
  \label{tab:scaling}
  \begin{tabular}{rcccccc}
    \toprule
    & \multicolumn{3}{c}{\textbf{Max ASU length before OOM} (residues)} & \multicolumn{3}{c}{\textbf{Wall-clock time at max ASU} (s)} \\
    \cmidrule(lr){2-4} \cmidrule(lr){5-7}
    $P$ & \textbf{1D} & \textbf{2D} & \textbf{Capacity ratio} & \textbf{1D} & \textbf{2D} & \textbf{Speedup} \\
    \midrule
    1  & 58  & 58  & $1.00\times$ & 1023 & 1025 & $1.00\times$ \\
    2  & 102 & --  & --           & 2340 & --   & --           \\
    3  & 120 & --  & --           & 2943 & --   & --           \\
    4  & 141 & 137 & $0.97\times$ & 4321 & 2230 & $1.94\times$ \\
    5  & 160 & --  & --           & 3845 & --   & --           \\
    6  & 173 & --  & --           & 3662 & --   & --           \\
    7  & 186 & --  & --           & 4802 & --   & --           \\
    8  & 196 & --  & --           & 4689 & --   & --           \\
    9  & 201 & 201 & $1.00\times$ & 7247 & 3253 & $2.23\times$ \\
    16 & 219 & 237 & $1.08\times$ & 7030 & 3785 & $1.86\times$ \\
    \bottomrule
  \end{tabular}
\end{table*}

\section{Scaling: max ASU size and inference time vs. number of GPUs}
\label{sec:results-scaling}

To validate the extent of applicability of our proposed CP techniques, 
we performed a scaling analysis on the task of designing proteins with icosahedral
point-group symmetry. Specifically, we target the icosahedral symmetry and sweep the 
amino-acid length of the asymmetric subunit (ASU) until inference runs out of memory (OOM)
for a fixed number of GPUs $P$. \Cref{tab:scaling} shows these values and the inference 
time for that specific run right before reaching OOM error. Unless explicitly mentioned,
all experiments in this subsection use HG200 GPUs with 95\,GB memory and a batch size of
one.

\subsection{Maximum ASU size.}
With either sharding scheme, the dominant quadratic pair tracks are distributed across 
the device mesh, so the memory ceiling is set primarily by the \emph{aggregate} available 
memory rather than by a single-device limit.  We believe this is the reason why we observe 
this value of max ASU length before OOM to be similar, independently of the sharding scheme
used.  Empirically, the maximum feasible ASU length increases with $P$ with the expected 
square-root trend predicted in \Cref{sec:cp-1d,sec:cp-2d}: as $P$ grows linearly while all the 
GPUs are maximally used, each device stores a quadratically smaller proportion of the 
growing pair tensors. 

\subsection{Inference time.}
While the two schemes reach comparable memory ceilings at a given $P$, their wall-clock 
scaling differs. The 1D scheme performs a single \textsc{AllGather} of the single-track 
activations after each attention block (\cref{sec:cp-1d}), which introduces a latency cost 
that grows with cluster size and impacts runtime. The 2D scheme replaces global gathers 
with $\sqrt{P}$-step ring communication over smaller shards and overlaps these shifts with 
local computation (\cref{sec:cp-2d}), yielding better time scaling in practice. This effect is 
reflected in \Cref{tab:scaling}, where 2D CP achieves lower per-step inference time at the 
same $P$, often finishing inference in around half the time with respect to the 1D inference.

% Ensure Table~\ref{tab:scaling} is placed before the next two-column figure.


\subsection{Designing octahedral nanoparticles on small GPUs}
\label{sec:results-democratization}

Octahedral protein cages are a promising but underexplored therapeutic scaffold class: their
rigid, multivalent geometry is suited to higher-order receptor engagement, and their 
interior volume can encapsulate biologic cargoes. Despite this potential, only a few 
\emph{de novo} octahedral cages have been characterised in therapy-relevant cell-based 
functional assays~\parencite{divine_designed_2021,yang_computational_2024}, and as noted in 
\cref{ch:current-memory-reduction-methods}, the memory cost of jointly sampling such large symmetric assemblies 
has been a practical barrier to broadening this design space.

To test whether Design-CP lifts this barrier on commodity hardware, we repeated the 
scaling protocol of \cref{sec:results-scaling} for octahedral targets, this time using a 
small cluster of workstation-grade NVIDIA RTX~A4000 GPUs ($16$\,GB per device,
${\sim}6\times$ less per-device memory than the HG200 GPUs used so far) and we report the 
results in \cref{app:O24_scaling_and_metrics}. In practice, the 2D sharding scheme allowed 
us to sample full octahedral assemblies up to $178$ residues per chain (with $24$ chains, 
this means modelling $4{,}272$ residues total) on $16$ A4000 GPUs 
(\cref{fig:octahedral}a), a per-chain size comparable to existing \emph{de novo} 
octahedral nanoparticles~\parencite{king_computational_2012}; we then evaluated $n=12$ such 
designs against the same families of \emph{in silico} metrics used for the icosahedral 
targets in \cref{sec:icosahedral}. Additional metrics and a discussion of the scaling 
sweep on this hardware are deferred to Appendix~\ref{app:O24_scaling_and_metrics}.

\subsection{Backbone sanity tracks the icosahedral baseline.} 
The backbone quality indicators (\cref{fig:octahedral}b--e) are healthy and
consistent with the icosahedral designs population: the per-chain count of chain breaks 
concentrates at zero with a single outlier at $2$, the count of backbone heavy-atom clashes 
is essentially zero across all designs (one outlier near $70$), the fraction of residues 
assigned to a helix or strand element sits at a median of $\approx 0.65$, and the maximum 
per-chain deviation of consecutive C$\alpha$--C$\alpha$ bond lengths from the canonical 
$3.8$\,\AA{} value clusters tightly around $\approx 0.3$\,\AA{}, well below the 
$0.75$\,\AA{} chain-break threshold. This last distribution is in fact noticeably tighter 
than for the icosahedral chains, and might suggest that the lower oligomeric state 
($24$ vs.\ $60$ subunits) imposes less geometric strain per chain.

\subsection{Symmetric interfaces are well-formed.} 
The two main symmetry-interface metrics shown in \cref{fig:octahedral}f--g are 
likewise healthy. The minimum C$\alpha$--C$\alpha$ distance between distinct chains sits 
in the $\approx 4$--$10$\,\AA{} band, consistent with proper inter-subunit packing without 
steric overlap, and the average number of C$\alpha$--C$\alpha$ contacts per inter-subunit 
interface centres around $\approx 30$ with a positive tail beyond $60$. An example of 
octahedral design is provided in \cref{fig:octahedral}a, where the chains organise 
into a closed, octahedrally symmetric cage. This confirms that the symmetrised trajectory 
realises the intended point-group geometry with favourable \emph{in silico} metrics on commodity 
GPUs.\\

Overall, these results indicate that large-assembly protein design can be made feasible 
on modest, widely available GPU hardware, lowering the barrier to entry for groups without 
access to large-memory accelerators.



\begin{figure}[t]
  \centering
  \includegraphics[width=\textwidth]{Octahedral_figure_main_text.png}
  \caption{\textbf{Designing octahedral nanoparticles on workstation-grade GPUs.}
   \textbf{a}, A representative Design-CP octahedral assembly with $176$ residues per 
   chain ($24$ chains, $4{,}224$ residues total) generated on $16$ NVIDIA RTX~A4000 GPUs 
   ($16$\,GB per device). \textbf{b--e}, Backbone-sanity metrics over $n=12$ such designs: 
   \textbf{b} chain breaks, \textbf{c} backbone clashes, \textbf{d} non-loop fraction, 
   \textbf{e} max CA deviation. \textbf{f--g}, Symmetry-interface metrics: \textbf{f} 
   minimum inter-chain distance, \textbf{g} mean contacts per interface. Metric definitions
    in Appendix~\ref{app:designability}; the remaining secondary-structure and contact 
    metrics are reported in Appendix~\ref{app:O24_scaling_and_metrics}.}
  \label{fig:octahedral}
\end{figure}


\section{Discussion}
\label{sec:discussion}

We introduced \emph{Design-CP}, two context-parallel inference strategies for RFDiffusion~3 
that shard the quadratic token- and atom-pair representations across multiple GPUs while 
preserving the model's architecture and weights. We characterised the memory ceilings and 
strong-scaling behaviour of a lightweight 1D row-sharding scheme and a 2D grid scheme 
(which requires a perfect-square GPU count $P$ to form a $\sqrt{P}\times\sqrt{P}$ mesh), 
and found that 2D sharding achieves better wall-clock scaling in practice by overlapping 
ring communication with computation.

Beyond RFD3, the underlying approach is model-agnostic: any protein design model whose 
inference is dominated by self-attention and other $\mathcal{O}(n^2)$ pairwise activations 
can, in principle, adopt the same sharding patterns. This suggests that context parallelism
is broadly applicable across modern design pipelines, which largely build on 
Transformer-style attention mechanisms with pari-bias.

We further showed that strong point-group symmetry constraints can make CP usable 
\emph{out of the box} for end-to-end all-atom design of large protein nanoparticles 
without retraining or fine-tuning. Prior work on context-parallelism for structure 
prediction (e.g., Fold-CP, \cite{lin_fold-cp_2026}) indicates that scaling inference to 
very large complexes can be limited by out-of-distribution generalisation once inputs exceed 
the model's native training regime. Our results refine this picture: for highly symmetric 
assemblies, symmetrised sampling collapses the effective degrees of freedom to a single ASU 
while still modelling the correct inter-subunit geometry, enabling CP to scale to full 
icosahedral and octahedral assemblies while preserving promising \emph{in silico} backbone-sanity 
and symmetry-interface metrics (\cref{sec:icosahedral,sec:results-democratization}). We 
also demonstrated that the same approach enables \emph{de novo}
design of octahedral nanoparticles on a small cluster of workstation-grade GPUs, indicating 
a practical path towards democratising large-assembly protein design.

A key limitation is that pushing beyond the native training crop can still introduce 
distribution shift and degrade outputs for highly asymmetric targets. Future work could 
therefore explore training or fine-tuning design models with longer contexts and/or explicit 
context-parallelism in the training loop. This is now not a priority for this project, as our
immediate goal is to design large symmetric assemblies for vaccine applications, but it is a 
natural next step for the broader design community.

% as well as experimental validation of the d
% esigned assemblies to assess how \emph{in silico} quality metrics translate to expression, 
% stability, and correct self-assembly \emph{in vitro}.
